\documentclass{ximera}
\title{Worksheet, Functions 1}

\newcommand{\pskip}{\vskip 0.1 in}

\begin{document}
\begin{abstract}
Introduction to functions.
\end{abstract}
\maketitle

\emph{Directions:} No calculators. Work with fractions.


\begin{question} \label{Qfvm334}
Complete the following sentences. Each - is a letter.

\begin{enumerate}
\item The domain of a relation is the  - - -  of its  - - - - - - .

\item The range of a relation is the - - - of its  - - - - - - - .
\end{enumerate}
\end{question}

\begin{question} \label{Q930f}
Find the domain and range of each function or relation.

\begin{enumerate}
\item $f(x) = \sqrt{25 - 49x}$

\item $g(x) = \frac{1}{\sqrt{25 - 49x}}$

\item $h(x) = \frac{31x - 13}{17 - 87x}$ (find just the domain for now)

\item $k(x) = 97 - \frac{2}{3}\left( x-12 \right)^2$

\item $\{ (x,y) \, | \, (x+8)^2 + (y-12)^2 = 5 \}$
\end{enumerate}
\end{question}

\begin{question} \label{Qor3or3}
Let 
\[
    f(x) = \frac{11}{7}\left( x - 9 \right) .
\]

\begin{enumerate}
\item Evaluate $f(-12)$.

\item Describe in complete sentences the action of $f$ on its input.

\item Use part (b) to solve the equation $f(x) = 55$.
\end{enumerate}
\end{question}

\begin{question} \label{Qffee}
Let 
\[
       g(x) = \frac{7}{3}\left(x-5\right)^2 +4
\]

\begin{enumerate}
\item Evaluate $g(-4)$.

\item Describe in complete sentences the action of $g$ on its input.

\item Use part (b) to solve the equation $g(x) = 18$.

\end{enumerate}
\end{question}

\begin{question} \label{Qdf3efr30}
Let $f(x) = x^2 + 12x + 3$ and solve the equation $f(x)=1$ \emph{without} using the quadratic formula.
\end{question}

\begin{question} \label{Q9emfe}
Let 
\[
    h(x) =12 -  (x+9)^2
\]

\begin{enumerate}
\item Evaluate $h(-13)$.

\item Describe the action of $h$ on its input.

\item Use part (b) to solve the equation $h(x) = -37$.

\item Solve the equation $h(x) = h(-5)$.

\item Solve the equation $h(x) = h(153)$.
\end{enumerate}

\begin{explanation}
\begin{enumerate}
\item
\begin{align*}
 h(-13) &= 12 - (-7+13)^2 \\
         &= 12 - (-6)^2 \\ 
         &= 12 - 36  \\
         &= -24
\end{align*}

\item \begin{itemize}
\item The function $h$ first adds $9$ to its input. 
\item It then squares the result.
\item It then multiplies that result by $-1$.
\item And finally it adds 12 to the previous result. 
\end{itemize}

\item To solve the equation $h(x) = -37$, we undo the actions of $h$ in the reverse order:

If $h(x) = -37$, then
\[
      12 -  (x+9)^2 = -37 .
\] 
And (subtracting 12 from both sides) this equation is true if and only if
\[
      - (x+9)^2 = -49,
\]
or (multiplying both sides by $-1$) if and only if
\[
 (x+9)^2 = 49.
\]
And since there are two numbers we can square to get $49$, this last equation is true if and only if
\[
  x+ 9 = 7 \hskip 0.2 in \text{or} \hskip 0.2 in  x+9 = -7
\]
\[
   x = -2  \hskip 0.2 in \text{or} \hskip 0.2 in  x = -16 .
\]
So the solution set to the equation $h(x)=-37$ is
\[
    \{x \, | \, x=-2 \text{ or } x=-16   \}.
\]


\item We start by evaluating $h(-5)$:
\begin{align*}
 h(-5) &= 12 - (-5+9)^2 \\
         &= 12 -4^2 \\ 
         &= 12 - 16  \\
         &= -4.
\end{align*}
The most common error is to stop here. But the question was not to evaluate $h(-5)$. Rather it was to solve the equation $h(x)=h(-5)$, which we can write as
\[
    12 -  (x+9)^2 = -4 .
\]
This is similar to part (c) and I'll leave the work to you. The final result is that the solution set to the equation $h(x) = h(-5)$ is
\[
   \{x \, | \, x=-5 \text{ or } x=-13   \}.
\]

\item We can do this without a calculator as follows.

The equation is
\[
   h(x) = h(153)
\]
or equivalently,
\[
  12 -  (x+9)^2 = 12 -  (153+9)^2 .
\]
And this equation is true if and only if
\[
   -  (x+9)^2 = -  (153+9)^2 ,
\]
or if and only if
\[
  (x+9)^2 = (153+9)^2.
\]
And this last equation is true if and only if
\[
   x+9 = 153+9 \hskip 0.2 in \text{or} \hskip 0.2 in x+9 = -(153+9)
\]
\[
      x = 153  \hskip 0.2 in \text{or} \hskip 0.2 in x = -171
\]
So the solution set to the equation $h(x) = h(153)$ is
\[
   \{x \, | \, x=153 \text{ or } x=-171  \}.
\]

\end{enumerate}


\end{explanation}

\end{question}

\begin{question} \label{Qef3er}
The function
\[
   g = f(v) , 25 \leq v \leq 75 ,
\]
expresses the gas mileage (in miles/gal) of a car in terms of its speed (in miles/hour). Translate each of the following into mathematics. Use \emph{complete} sentences.

\begin{enumerate}

\item What is the car's gas mileage at a speed of 40 miles/hour?


\emph{Solution:} To determine the car's gas mileage at a speed of $40$ miles/hour, we would ....

\item At what speeds does the car get 32 miles/gal?


\emph{Solution:} To determine the speeds at which the car gets $32$ miles/gal we would ...


\item At what speeds does the car get at most 32 miles/gal?


\emph{Solution:} To determine the speed at which the car gets at most $32$ miles/gallon, we would ...

\item At what speeds does the car get the same gas mileage as it does at a speed of 37 miles/hour?

\item At what speeds is the gas mileage $3$ miles/gal less than it is at a speed of 30 miles/hour?

\item At what speeds is the gas mileage at least $5$ miles/gal greater than it is at a speed of $30$ miles/hour?

\item At what speeds is the gas mileage $25\%$ less than it is at a speed of $50$ miles/hour?

\item At what speeds does decreasing the speed by $7$ miles/hour leave the gas mileage unchanged?

\item At what speeds does increasing the speed by $5$ miles/hour decrease the gas mileage by $2$ miles/gal?

\item At what speeds does increasing the speed by $5$ miles/hour increase the gas mileage by $10\%$?

\end{enumerate}
\end{question}

\begin{question} \label{Q9efr3r}
The function
\[
    g = f(v) = 40-\frac{1}{25}\left(v-50\right)^{2} \, , \, 25\leq v \leq 70 ,
\]
expresses the gas mileage (in miles/gal) of a car in terms of its speed (in miles/hr).

Use this function to answer all parts of Question 7.
\end{question}


\end{document}


